\begin{abstract}
	The most widely used \ransac variants score a candidate model by summing
	per-point scores that fall to zero beyond a threshold, and that threshold
	must be tuned to the noise level of the inliers. Recent work removed it by
	reversing the usual order of inference: for each split of the data into
	inliers and outliers, the noise level is integrated out in closed form, and
	the best split is then chosen. One setting with units of distance remained:
	the width of the region over which the outliers are spread. We remove it as
	well. Choosing the best split also tells us, for every candidate model,
	which points are outliers, and the farthest of them show how far the
	outliers reach. We estimate this reach with extreme value theory, by
	fitting the distribution of the outlier distances that exceed a high
	quantile. The resulting score is the first \ransac score that needs no
	threshold, no inlier noise level and no outlier width: changing the units of
	the residuals does not change which model it selects. Its one remaining
	setting is a quantile, which has no units and to which the score is nearly
	insensitive. Unlike earlier uses of extreme value theory in \ransac, which
	model descriptor matching scores to guide sampling, we model the geometric
	residuals of each candidate to score it. \TODO{one sentence of real-data
	results}
\end{abstract}
